% GENERATED FILE. Edit canonical lessons, then run npm run generate:books.
% canonical-source-hash: fnv1a64:73357b3fced11c94
% canonical-lessons: LA-C384-suffragium, LA-C384-comitia, LA-C384-candidatus, LA-C384-ambitio, LA-C384-eligo

\chapter{Elections}
\label{ch:la-elections}
\hlchaptermodality{384}

\begin{chapteropening}
\textbf{By the end of this chapter:} \emph{I can talk about votes and the voting assembly, a candidate, canvassing, and choosing someone.}

Complete the last lesson of chapter 384: I can talk about votes and the voting assembly, a candidate, canvassing, and choosing someone.
\end{chapteropening}

\section[suffrāgium, suffrāgiī]{suffrāgium, suffrāgiī — a vote; the right to vote}

\label{lesson:LA-C384-suffragium}



\begin{quote}
Before the new one: say the Latin for a quaestor, then the Latin for a censor.
\end{quote}

\subsection*{You'll want to know: suffrāgium, suffrāgiī}
\textbf{suffrāgium, suffrāgiī} — \textquotedblleft{}a vote; the right to vote\textquotedblright{}.

\begin{quote}
Cīvēs suffrāgia ferunt.
\end{quote}

\begin{quote}
\emph{The citizens cast their votes.}
\end{quote}

\emph{suffrāgium, suffrāgiī} n. \emph{Suffrāgium ferre}, to vote. English \emph{suffrage}.

\begin{cousinweb}
The origin of \emph{suffrāgium} is not settled: Lewis \& Short link it, with a \textquotedblleft{}perhaps\textquotedblright{}, to \emph{suffrāgō} \textquotedblleft{}the hock of a leg\textquotedblright{}, and EtymDB has no parent for it. English \emph{suffrage} comes from it.
\end{cousinweb}

\begin{grammarlens}[title={What you've built}]
One more word for telling, explaining and describing at B1.
\end{grammarlens}

\subsection*{Your turn}
\begin{itemize}
\raggedright
  \item \emph{Say it:} \emph{suffrāgium, suffrāgiī}
  \item \emph{Say it:} \emph{suffrāgium, suffrāgiī}, once more
  \item \emph{Say it:} a sentence of your own with \emph{suffrāgium, suffrāgiī}
  \item \emph{Recall:} say the Latin for a quaestor, then the Latin for a censor, then say \emph{suffrāgium, suffrāgiī} again
\end{itemize}

\begin{tcolorbox}[breakable,colback=teal!4,colframe=teal!35!black,title={Before you move on}]
What does suffrāgium, suffrāgiī mean? (\textquotedblleft{}A vote; the right to vote\textquotedblright{}.) Say it once more.
\end{tcolorbox}
\section[comitia, comitiōrum]{comitia, comitiōrum — the voting assembly, the elections}

\label{lesson:LA-C384-comitia}



\begin{quote}
Before the new one: say the Latin for a censor, then the Latin for a vote.
\end{quote}

\subsection*{You'll want to know: comitia, comitiōrum}
\textbf{comitia, comitiōrum} — \textquotedblleft{}the voting assembly, the elections\textquotedblright{}.

\begin{quote}
Comitia mox habēbuntur.
\end{quote}

\begin{quote}
\emph{The elections will be held soon.}
\end{quote}

\emph{comitia, comitiōrum} n. pl.: the Roman people meeting to vote. The singular \emph{comitium} is the meeting place.

\begin{cousinweb}
\emph{comitia} is built from \emph{com-} \textquotedblleft{}together\textquotedblright{} and the root of \emph{eō} \textquotedblleft{}to go\textquotedblright{} (Lewis \& Short): a going together.
\end{cousinweb}

\begin{grammarlens}[title={What you've built}]
One more word for telling, explaining and describing at B1.
\end{grammarlens}

\subsection*{Your turn}
\begin{itemize}
\raggedright
  \item \emph{Say it:} \emph{comitia, comitiōrum}
  \item \emph{Say it:} \emph{comitia, comitiōrum}, once more
  \item \emph{Say it:} a sentence of your own with \emph{comitia, comitiōrum}
  \item \emph{Recall:} say the Latin for a censor, then the Latin for a vote, then say \emph{comitia, comitiōrum} again
\end{itemize}

\begin{tcolorbox}[breakable,colback=teal!4,colframe=teal!35!black,title={Before you move on}]
What does comitia, comitiōrum mean? (\textquotedblleft{}The voting assembly, the elections\textquotedblright{}.) Say it once more.
\end{tcolorbox}
\section[candidātus, candidātī]{candidātus, candidātī — a candidate}

\label{lesson:LA-C384-candidatus}



\begin{quote}
Before the new one: say the Latin for a vote, then the Latin for the voting assembly.
\end{quote}

\subsection*{You'll want to know: candidātus, candidātī}
\textbf{candidātus, candidātī} — \textquotedblleft{}a candidate\textquotedblright{}.

\begin{quote}
Candidātus cīvēs in forō salūtābat.
\end{quote}

\begin{quote}
\emph{The candidate was greeting the citizens in the Forum.}
\end{quote}

\emph{candidātus, candidātī} m., from \emph{candidus}, shining white: a man seeking office wore a toga whitened with chalk. English \emph{candidate}.

\begin{cousinweb}
\emph{candidātus} is built on \emph{candidus} \textquotedblleft{}shining white\textquotedblright{} (EtymDB; Lewis \& Short): dressed in white. English \emph{candidate} comes from it.
\end{cousinweb}

\begin{grammarlens}[title={What you've built}]
One more word for telling, explaining and describing at B1.
\end{grammarlens}

\subsection*{Your turn}
\begin{itemize}
\raggedright
  \item \emph{Say it:} \emph{candidātus, candidātī}
  \item \emph{Say it:} \emph{candidātus, candidātī}, once more
  \item \emph{Say it:} a sentence of your own with \emph{candidātus, candidātī}
  \item \emph{Recall:} say the Latin for a vote, then the Latin for the voting assembly, then say \emph{candidātus, candidātī} again
\end{itemize}

\begin{tcolorbox}[breakable,colback=teal!4,colframe=teal!35!black,title={Before you move on}]
What does candidātus, candidātī mean? (\textquotedblleft{}A candidate\textquotedblright{}.) Say it once more.
\end{tcolorbox}
\section[ambitiō, ambitiōnis]{ambitiō, ambitiōnis — canvassing; ambition}

\label{lesson:LA-C384-ambitio}



\begin{quote}
Before the new one: say the Latin for the voting assembly, then the Latin for a candidate.
\end{quote}

\subsection*{You'll want to know: ambitiō, ambitiōnis}
\textbf{ambitiō, ambitiōnis} — \textquotedblleft{}canvassing; ambition\textquotedblright{}.

\begin{quote}
Ambitiō eius omnibus nōta erat.
\end{quote}

\begin{quote}
\emph{His canvassing was known to all.}
\end{quote}

\emph{ambitiō, ambitiōnis} f., built on \emph{ambiō}, to go round: the candidate went round asking for votes. English \emph{ambition}.

\begin{cousinweb}
\emph{ambitiō} is the noun built on \emph{ambiō} \textquotedblleft{}to go round\textquotedblright{} (EtymDB; Lewis \& Short), from \emph{amb-} \textquotedblleft{}around\textquotedblright{} and \emph{eō} \textquotedblleft{}to go\textquotedblright{}. English \emph{ambition} comes from it.
\end{cousinweb}

\begin{grammarlens}[title={What you've built}]
One more word for telling, explaining and describing at B1.
\end{grammarlens}

\subsection*{Your turn}
\begin{itemize}
\raggedright
  \item \emph{Say it:} \emph{ambitiō, ambitiōnis}
  \item \emph{Say it:} \emph{ambitiō, ambitiōnis}, once more
  \item \emph{Say it:} a sentence of your own with \emph{ambitiō, ambitiōnis}
  \item \emph{Recall:} say the Latin for the voting assembly, then the Latin for a candidate, then say \emph{ambitiō, ambitiōnis} again
\end{itemize}

\begin{tcolorbox}[breakable,colback=teal!4,colframe=teal!35!black,title={Before you move on}]
What does ambitiō, ambitiōnis mean? (\textquotedblleft{}Canvassing; ambition\textquotedblright{}.) Say it once more.
\end{tcolorbox}
\section[ēligō, ēligere]{ēligō, ēligere — to pick out, to choose}

\label{lesson:LA-C384-eligo}



\begin{quote}
Before the new one: say the Latin for a candidate, then the Latin for canvassing.
\end{quote}

\subsection*{You'll want to know: ēligō, ēligere}
\textbf{ēligō, ēligere} — \textquotedblleft{}to pick out, to choose\textquotedblright{}.

\begin{quote}
Ex multīs hunc ēlēgimus.
\end{quote}

\begin{quote}
\emph{Out of many we chose this one.}
\end{quote}

\emph{ēligō, ēligere, ēlēgī, ēlēctum}: \emph{ē} + \emph{legō}, to pick. English \emph{elect} and \emph{election}.

\begin{cousinweb}
\emph{ēligō} is built from \emph{ē} \textquotedblleft{}out\textquotedblright{} and \emph{legō} \textquotedblleft{}to pick, to gather\textquotedblright{} (EtymDB; Lewis \& Short). English \emph{elect} comes from its participle \emph{ēlēctus}.
\end{cousinweb}

\begin{grammarlens}[title={What you've built}]
One more word for telling, explaining and describing at B1.
\end{grammarlens}

\subsection*{Your turn}
\begin{itemize}
\raggedright
  \item \emph{Say it:} \emph{ēligō, ēligere}
  \item \emph{Say it:} \emph{ēligō, ēligere}, once more
  \item \emph{Say it:} a sentence of your own with \emph{ēligō, ēligere}
  \item \emph{Recall:} say the Latin for a candidate, then the Latin for canvassing, then say \emph{ēligō, ēligere} again
\end{itemize}

\begin{tcolorbox}[breakable,colback=teal!4,colframe=teal!35!black,title={Before you move on}]
What does ēligō, ēligere mean? (\textquotedblleft{}To pick out, to choose\textquotedblright{}.) Say it once more.
\end{tcolorbox}
